\section{Глава 5. Методы Монте-Карло}
\subsection{Упражнение 5.1(c. 126)}
Рассмотрим диаграммы на рис. $5.1$ справа. Почему оценка функции ценности резко возрастает в последних двух строках на заднем плане? Почему она спадает в последней строке на диаграмме слева? Почему значения на переднем плане на верхних диаграммах выше, чем на нижних? 
\begin{figure}[H]
    \centering
    \includegraphics[width=1\linewidth]{fig_5_1.png}
    \caption{Рис. 5.1. Приближенные функции ценности состояний для стратегии игры в блэкджек, при которой игрок останавливается, набрав $20$ или $21$. Вычислены посредством оценивания стратегии методом Монте-Карло.}
\end{figure}
\begin{solution}
    \begin{enumerate}
        \item Из условия задачи очевидно, что стратегия игрока далеко не оптимальна: дойти до суммы $20$ или $21$ сложно без "перебора". Но в случае достижения таких значений выигрыш практически гарантирован, поэтому ценность состояний резко возрастает на последних двух строках.
        \item Левая строка соответствует случаю, когда открытой картой сдающего является туз. Это самый выгодный случай для противника, поэтому ценность состояний падает в этой строке.
        \item Ценность состояний при наличии играющего туза больше ценности состояний в случае его отсутствия, потому что играющий туз дает игроку запасной шанс: при переборе можно засчитать туз за $1$ и продолжить игру.
    \end{enumerate}
\end{solution}

\subsection{Упражнение 5.2(c. 126)}
Предположим, что вместо МК первого посещения для моделирования игры в блэкджек использовался МК всех посещений. Как вы думаете, будет ли результат сильно отличаться? Объясните свой ответ. 
\begin{solution}
    Не будет, так как внутри одного эпизода состояния повторяться не могут(так как сумма очков у игрока может только увеличиться). Это означает, что первое посещение состояния будет единственным. Небольшие различия могут быть обусловлены шумом.
\end{solution}

\subsection{Упражнение 5.3(c. 128)}
Нарисуйте диаграмму предшествующих состояний для оценивания $q_{\pi}$ методом Монте-Карло. 
\begin{solution}
    \leavevmode

    \medskip
    \centering
    \includegraphics[width=0.4\linewidth]{fig_5_3_sol.png}

    \captionof{figure}{
    По аналогии с функциями ценности состояний теперь мы оцениваем пару состояние-действие, которая зависит от следующих пар до конца эпизода. Ветвления нет, потому что мы считаем доход по эпизоду.
    }
    \label{fig:5_3_sol}
\end{solution}

\subsection{Упражнение 5.4(c. 131)}
Псевдокод метода Монте-Карло ИС не эффективен, потому что для каждой пары состояние–действие он хранит список всех доходов и повторно вычисляет среднее. Было бы эффективнее использовать описанную в разделе 2.4
технику – хранить только среднее и количество элементов в списке (для каждой пары состояние–действие) и обновлять их инкрементно. Покажите, как для этого следует модифицировать псевдокод.
\begin{solution}
    \leavevmode\\
    \begin{algorithm}[H]
    \SetAlgoLined
    \DontPrintSemicolon
    \caption{Метод Монте-Карло ИС (с исследовательскими стартами), инкрементная версия}

    \textbf{Инициализация:}\;
    \Indp
    $\pi(s) \in \mathcal{A}(s)$ (произвольная) для всех $s \in \mathcal{S}$\;
    $Q(s, a) \in \mathbb{R}$ (произвольная) для всех $s \in \mathcal{S},\ a \in \mathcal{A}(s)$\;
    $N(s, a) \leftarrow 0$ для всех $s \in \mathcal{S},\ a \in \mathcal{A}(s)$\;
    \Indm

    \BlankLine
    \textbf{Повторять бесконечно (для каждого эпизода):}\;
    \Indp
    Выбрать $S_0 \in \mathcal{S},\ A_0 \in \mathcal{A}(S_0)$ случайным образом, так что вероятность любой пары больше 0\;
    Сгенерировать эпизод, начинающийся в $S_0, A_0$ и следующий $\pi$:\;
    \quad $S_0, A_0, R_1, \ldots, S_{T-1}, A_{T-1}, R_T$\;
    $G \leftarrow 0$\;
    \For{$t = T-1, T-2, \ldots, 0$}{
        $G \leftarrow \gamma G + R_{t+1}$\;
        \If{пара $S_t, A_t$ не встречается в $S_0, A_0, S_1, A_1, \ldots, S_{t-1}, A_{t-1}$}{
            $N(S_t, A_t) \leftarrow N(S_t, A_t) + 1$\;
            $Q(S_t, A_t) \leftarrow Q(S_t, A_t) + \dfrac{1}{N(S_t, A_t)}\bigl(G - Q(S_t, A_t)\bigr)$\;
            $\pi(S_t) \leftarrow \arg\max_a Q(S_t, a)$\;
        }
    }
    \Indm

    \end{algorithm}
\end{solution}

\subsection{Упражнение 5.5(c. 138)}
Рассмотрим МППР с одним незаключительным состоянием и одним действием, которое ведет назад в незаключительное состояние с вероятностью $p$ и в заключительное состояние с вероятностью $1 - p$. Пусть вознаграждение
за любой переход равно $+1$ и $\gamma = 1$. Предположим, что наблюдается один эпизод,
продолжающийся $10$ шагов, с доходом $10$. Каковы оценки ценности незаключительного состояния в случае первого посещения и всех посещений? 
\begin{solution}
    Эпизод длиной $10$ шагов с доходом $G = 10$ имеет вид:

    \begin{center}
    \resizebox{\textwidth}{!}{%
    \begin{tikzpicture}[
        every node/.style={font=\small},
        state/.style={circle, draw, thick, minimum size=7mm, inner sep=1pt},
        terminal/.style={rectangle, draw, thick, minimum size=7mm, inner sep=2pt, fill=gray!20},
        arrow/.style={-{Stealth[length=2mm]}, thick}
    ]
    \foreach \i in {0,...,9} {
        \node[state] (s\i) at (\i*1.4, 0) {$s$};
        \node[below=1mm of s\i] {\scriptsize $t=\i$};
    }
    \node[terminal] (term) at (14, 0) {\scriptsize term};
    \node[below=1mm of term] {\scriptsize $t=10$};
    \foreach \i [evaluate=\i as \j using int(\i+1)] in {0,...,8} {
        \draw[arrow] (s\i) -- node[above, font=\scriptsize] {$+1$} (s\j);
    }
    \draw[arrow] (s9) -- node[above, font=\scriptsize] {$+1$} (term);
    \end{tikzpicture}%
    }
    \end{center}

    При $\gamma = 1$ доход с момента $t$ равен $G_t = 10 - t$.

    \textbf{First-visit MC:} учитывается только первое посещение ($t = 0$):
    \[
    V_{\text{first-visit}}(s) = G_0 = 10.
    \]

    \textbf{Every-visit MC:} усредняются доходы со всех $10$ посещений:
    \[
    V_{\text{every-visit}}(s) = \frac{10 + 9 + \cdots + 1}{10} = \frac{55}{10} = 5{,}5.
    \]
\end{solution}

\subsection{Упражнение 5.6(c. 141)}
Как выглядит уравнение, аналогичное $(5.6)$, если вместо ценностей состояний $V(s)$ использовать ценности действий $Q(s, a)$, опять-таки при известных доходах, сгенерированных с помощью стратегии $b$?
\begin{solution}
    В формуле (5.6) для оценки $V(s)$:
    \[
    V(s) \doteq \frac{\sum_{t \in \mathcal{T}(s)} \rho_{t:T(t)-1}\, G_t}{\sum_{t \in \mathcal{T}(s)} \rho_{t:T(t)-1}}
    \]
    нужно сделать две замены.

    \textbf{1. Множество индексов.} Вместо $\mathcal{T}(s)$ — моментов посещения состояния $s$ — берём
    \[
    \mathcal{T}(s,a) = \{t : S_t = s,\ A_t = a\}.
    \]

    \textbf{2. Коэффициент $\rho$ начинается с шага $t+1$.} Действие $A_t = a$
    зафиксировано условием, поэтому отношение $\pi(A_t \mid S_t) / b(A_t \mid S_t)$
    для шага $t$ в коэффициент не входит — корректировать нужно только то, что
    происходит после этого выбора:
    \[
    \rho_{t+1:T(t)-1} = \prod_{k=t+1}^{T(t)-1} \frac{\pi(A_k \mid S_k)}{b(A_k \mid S_k)}.
    \]

    \textbf{Итоговая формула:}
    \[
    \boxed{\,Q(s,a) \doteq \frac{\sum_{t \in \mathcal{T}(s,a)} \rho_{t+1:T(t)-1}\, G_t}{\sum_{t \in \mathcal{T}(s,a)} \rho_{t+1:T(t)-1}}\,}
    \]
\end{solution}

\subsection{Упражнение 5.7(c. 141)}
На кривых обучения типа показанных на рис. $5.3$ ошибка обычно убывает по мере обучения, как и случилось для метода обыкновенной выборки по значимости. Но для метода взвешенной выборки по значимости ошибка сначала
возрастает, а потом убывает. Как вы думаете, почему так происходит?
\begin{figure}[H]
    \centering
    \includegraphics[width=1\linewidth]{fig_5_3.png}
    \caption{Рис. 5.3. Взвешенная выборка по значимости дает меньшую ошибку оценки ценности одного состояния в блэкджеке по эпизодам с разделенной стратегией}
\end{figure}
\begin{solution}
    Оценка инициализируется нулём, что близко к истинной ценности $-0{,}27726$,
    поэтому начальная ошибка мала. Поведенческая стратегия выбирает «ещё» и
    «хватит» равновероятно, тогда как целевая всегда останавливается, поэтому
    эпизоды с ненулевым $\rho$ редки.

    Пока все наблюдённые $\rho_t = 0$, знаменатель нулевой и оценка остаётся
    равной начальному значению $0$. Как только появляется первый эпизод с
    $\rho > 0$, веса в формуле
    \[
    V(s) = \frac{\sum_t \rho_t\, G_t}{\sum_t \rho_t}
    \]
    сокращаются, и оценка становится равна доходу этого единственного эпизода:
    \[
    V(s) = \frac{\rho_1 G_1}{\rho_1} = G_1 \in \{-1, 0, +1\}.
    \]
    Прыжок от $0$ к $\pm 1$ резко увеличивает квадратичную ошибку относительно
    истинного значения $-0{,}28$. По мере накопления согласующихся эпизодов
    оценка усредняется и подтягивается к истинной ценности — ошибка убывает.

    Иными словами, всплеск возникает потому, что взвешенная оценка по одному
    эпизоду полностью игнорирует вес и сводится к самому доходу, что временно
    уводит её от истинной ценности.
\end{solution}

\subsection{Упражнение 5.8(c. 141)}
При получении результатов в примере $5.5$, показанных на рис. $5.4$,
использовался метод МК первого посещения. Допустим, что для той же задачи был
применен метод МК всех посещений. Будет ли в таком случае дисперсия оценки
по-прежнему бесконечной? Объясните свой ответ.
\begin{figure}[H]
    \centering
    \includegraphics[width=1\linewidth]{fig_5_4.png}
    \caption{Рис. 5.4. Обыкновенная выборка по значимости порождает на удивление неустойчивые оценки для МППР с  одним состоянием, показанным
    в верхней части (пример 5.5). В данном случае правильная оценка равна 1
    $(\gamma = 1)$, но, хотя это действительно математическое ожидание выборочного дохода (после выборки по значимости), дисперсия выборки бесконечна
    и оценки не сходятся к этому значению. Эти результаты имеют место для МК
    первого посещения с разделенной стратегией
    }
\end{figure}
\begin{solution}
    Дисперсия по-прежнему бесконечна.

    В методе всех посещений каждый эпизод даёт вклад в $\mathbb{E}_b[\rho^2 G^2]$
    от \emph{каждого} посещения состояния $s$, тогда как в методе первого
    посещения учитывается только посещение в момент $t = 0$. Слагаемое от
    первого посещения в обоих методах одно и то же, а в методе всех посещений
    к нему добавляются дополнительные слагаемые от посещений в моменты
    $t = 1, 2, \ldots$

    Все слагаемые неотрицательны ($\rho^2 \geq 0$, $G^2 \geq 0$), поэтому
    $\mathbb{E}_b[\rho^2 G^2]$ для метода всех посещений не меньше, чем для
    метода первого посещения. Поскольку в учебнике для first-visit получен
    расходящийся ряд $0{,}2\sum_{k=0}^\infty 1{,}8^k = \infty$, больший по
    членам ряд для every-visit расходится тем более — дисперсия остаётся
    бесконечной.
\end{solution}

\subsection{Упражнение 5.9(c. 142)}
Модифицируйте алгоритм для оценивания стратегии методом МК первого посещения (раздел $5.1$), воспользовавшись инкрементной реализацией выборочного среднего, описанной в разделе $2.4$.
\begin{solution}
    \leavevmode\\
    \begin{algorithm}[H]
    \SetAlgoLined
    \DontPrintSemicolon
    \caption{Предсказание методом МК первого посещения для оценивания $V \approx v_\pi$ (инкрементная версия)}

    \KwIn{стратегия $\pi$, подлежащая оцениванию}

    \textbf{Инициализация:}\;
    \Indp
    $V(s) \in \mathbb{R}$ (произвольная) для всех $s \in \mathcal{S}$\;
    $N(s) \leftarrow 0$ для всех $s \in \mathcal{S}$\;
    \Indm

    \BlankLine
    \textbf{Повторять бесконечно (для каждого эпизода):}\;
    \Indp
    Сгенерировать эпизод, следуя $\pi$: $S_0, A_0, R_1, S_1, A_1, R_2, \ldots, S_{T-1}, A_{T-1}, R_T$\;
    $G \leftarrow 0$\;
    \For{$t = T-1, T-2, \ldots, 0$}{
        $G \leftarrow \gamma G + R_{t+1}$\;
        \If{$S_t$ не встречается в $S_0, S_1, \ldots, S_{t-1}$}{
            $N(S_t) \leftarrow N(S_t) + 1$\;
            $V(S_t) \leftarrow V(S_t) + \dfrac{1}{N(S_t)}\bigl(G - V(S_t)\bigr)$\;
        }
    }
    \Indm

    \end{algorithm}
\end{solution}

\subsection{Упражнение 5.10(c. 143)}
Выведите правило обновления взвешенного среднего $(5.8)$ из
$(5.7)$. Рассуждайте так же, как при выводе правила для обычного среднего $(2.3)$.
\begin{equation}
    V_n \doteq \frac{\sum_{k=1}^{n-1} W_k G_k}{\sum_{k=1}^{n-1} W_k}, \quad n \geq 2,
    \tag{5.7}
\end{equation}
\begin{equation}
    V_{n+1} \doteq V_n + \frac{W_n}{C_n}\bigl[G_n - V_n\bigr], \quad n \geq 1,
    \tag{5.8}
\end{equation}
\begin{solution}
    Обозначим $C_n = \sum_{k=1}^{n} W_k$. Тогда из $(5.7)$:
    \[
    V_{n+1} = \frac{\sum_{k=1}^{n} W_k G_k}{C_n} = \frac{\sum_{k=1}^{n-1} W_k G_k + W_n G_n}{C_n}.
    \]
    Заметим, что $\sum_{k=1}^{n-1} W_k G_k = C_{n-1} V_n$, поэтому
    \[
    V_{n+1} = \frac{C_{n-1} V_n + W_n G_n}{C_n} = \frac{(C_n - W_n) V_n + W_n G_n}{C_n} = V_n + \frac{W_n}{C_n}\bigl[G_n - V_n\bigr].
    \]
\end{solution}

\subsection{Упражнение 5.11(c. 144)}
В приведённом выше алгоритме управления методом МК с разделённой стратегией вы, возможно, ожидали, что в обновлении $W$ будет участвовать коэффициент выборки по значимости $\pi(A_t \mid S_t) / b(A_t \mid S_t)$, а на самом деле мы видим $1/b(A_t \mid S_t)$. Почему же, несмотря на это, алгоритм корректен?
\begin{solution}
    Целевая стратегия $\pi$ детерминирована и жадна: $\pi(s) = \arg\max_a Q(s, a)$. Поэтому $\pi(A_t \mid S_t)$ принимает лишь два значения:
    \[
    \pi(A_t \mid S_t) = 
    \begin{cases}
        1, & \text{если } A_t = \pi(S_t), \\
        0, & \text{если } A_t \neq \pi(S_t).
    \end{cases}
    \]
    В случае $A_t \neq \pi(S_t)$ весь коэффициент $\rho$ обнуляется, и вклад в обновление становится нулевым. В оставшемся случае $A_t = \pi(S_t)$ выполнено $\pi(A_t \mid S_t) = 1$, поэтому
    \[
    \frac{\pi(A_t \mid S_t)}{b(A_t \mid S_t)} = \frac{1}{b(A_t \mid S_t)}.
    \]
\end{solution}

\subsection{Упражнение 5.12. Кольцевые гонки(c. 145)}
Рассмотрим
управление гоночным автомобилем на повороте, показанном на рис.  $5.5$. Вы
хотите пройти его как можно быстрее, но не настолько быстро, чтобы вылететь
с трассы. В нашей упрощенной модели кольцевых гонок автомобиль может находиться в одной из дискретных ячеек сетки. Скорость также дискретна и равна
количеству ячеек, на которые машина переместилась по горизонтали и по вертикали за один временной шаг. Действия – приращения составляющих скорости.
Каждая может изменяться на $+1, -1$ или $0$ на одном шаге, так что всего имеется
девять $(3 \times 3)$ действий. Обе составляющие скорости должны быть неотрицательны и меньше $5$, и одновременно могут быть равны $0$ только на линии старта. Каждый эпизод начинается в одном из случайно выбранных начальных состояний,
когда обе составляющие скорости равны нулю, а  заканчивается, когда машина
пересечет линию финиша. За каждый шаг до пересечения линии финиша полагается вознаграждение $-1$. Если машина задела границу трассы, то она возвращается
в случайную позицию на линии старта, обе составляющие скорости обнуляются,
и  эпизод продолжается. Прежде чем обновлять положение машины на каждом
шаге, следует проверить, не пересекает ли предполагаемая траектория движения
границу трассы. Если машина пересекает линию финиша, то эпизод заканчивается; если она пересекает какую-то другую границу, то считается, что машина вылетела с трассы, поэтому она возвращается на линию старта. Чтобы сделать задачу
более интересной, с вероятностью $0.1$ на каждом временном шаге оба приращения, составляющих скорости, равны нулю независимо от того, какими они должны
были бы стать. Примените к этой задаче метод управления Монте-Карло, чтобы
вычислить оптимальную стратегию при старте из каждого начального состояния.
Нарисуйте несколько траекторий, следующих оптимальной стратегии (но подавите для них шум).
\begin{figure}[H]
    \centering
    \includegraphics[width=0.6\linewidth]{fig_5_5.png}
    \caption{Рис. 5.5. Два правых поворота в задаче о круговых гонках}
\end{figure}
\begin{solution}
    \textbf{Код смотри в отдельном файле репозитория.} \\
    \begin{figure}[H]
    \centering
    \begin{minipage}{0.4\textwidth}
        \centering
        \includegraphics[width=\linewidth]{track_a.png}
    \end{minipage}\hfill
    \begin{minipage}{0.4\textwidth}
        \centering
        \includegraphics[width=\linewidth]{track_b.png}
    \end{minipage}
    \end{figure}
\end{solution}


\subsection{*Упражнение 5.13(c. 148)}
Восполните шаги, необходимые для вывода $(5.14)$ из $(5.12)$.
\begin{equation}
    \resizebox{0.8\textwidth}{!}{$\displaystyle \rho_{t:T-1} R_{t+1} = \frac{\pi(A_t|S_t)}{b(A_t|S_t)} \frac{\pi(A_{t+1}|S_{t+1})}{b(A_{t+1}|S_{t+1})} \frac{\pi(A_{t+2}|S_{t+2})}{b(A_{t+2}|S_{t+2})} \cdots \frac{\pi(A_{T-1}|S_{T-1})}{b(A_{T-1}|S_{T-1})} R_{t+1}. $}
    \tag{5.12}
\end{equation}
\begin{equation}
    \mathbb{E}[\rho_{t:T-1} R_{t+1}] = \mathbb{E}[\rho_{t:t} R_{t+1}].
    \tag{5.14}
\end{equation}
\begin{solution}
    Подставляя $(5.12)$ в левую часть $(5.14)$, получаем:
    \[
    \mathbb{E}[\rho_{t:T-1} R_{t+1}] = \mathbb{E}\left[\frac{\pi(A_t|S_t)}{b(A_t|S_t)} R_{t+1} \cdot \prod_{k=t+1}^{T-1} \frac{\pi(A_k|S_k)}{b(A_k|S_k)}\right].
    \]
    Матожидание произведения распадается на произведение матожиданий:
    \[
    = \mathbb{E}[\rho_{t:t} R_{t+1}] \cdot \prod_{k=t+1}^{T-1} \mathbb{E}\left[\frac{\pi(A_k|S_k)}{b(A_k|S_k)}\right].
    \]
    По $(5.13)$ каждый множитель при $k \geq t+1$ равен $1$, откуда
    \[
    \mathbb{E}[\rho_{t:T-1} R_{t+1}] = \mathbb{E}[\rho_{t:t} R_{t+1}]. 
    \]
\end{solution}

\subsection{*Упражнение 5.14(c. 148)}
Модифицируйте алгоритм управления Монте-Карло с разделённой стратегией (стр.~144), воспользовавшись идеей усечённой средневзвешенной оценки $(5.10)$. Заметим, что сначала нужно будет преобразовать это равенство, введя ценности действий.
\begin{equation}
    \resizebox{0.85\linewidth}{!}{$\displaystyle V(s) \doteq \frac{\sum_{t \in \mathcal{T}(s)} \left((1-\gamma)\sum_{h=t+1}^{T(t)-1} \gamma^{h-t-1} \rho_{t:h-1} \bar{G}_{t:h} + \gamma^{T(t)-t-1} \rho_{t:T(t)-1} \bar{G}_{t:T(t)}\right)}{\sum_{t \in \mathcal{T}(s)} \left((1-\gamma)\sum_{h=t+1}^{T(t)-1} \gamma^{h-t-1} \rho_{t:h-1} + \gamma^{T(t)-t-1} \rho_{t:T(t)-1}\right)}. $}
    \tag{5.10}
\end{equation}
\begin{solution}
    \textbf{Преобразование $(5.10)$ для ценностей действий.}
    Действие $A_t = a$ задано условием, поэтому importance ratio начинаются с $\rho_{t+1:\cdot}$. Выделим слагаемое $h = t+1$ из внутренней суммы (для него $\bar{G}_{t:t+1} = R_{t+1}$ и importance ratio отсутствует):
    \[
    \resizebox{0.97\linewidth}{!}{$\displaystyle Q(s,a) \doteq \frac{\sum_{t \in \mathcal{T}(s,a)} \left((1-\gamma) R_{t+1} + (1-\gamma)\sum_{h=t+2}^{T(t)-1} \gamma^{h-t-1} \rho_{t+1:h-1} \bar{G}_{t:h} + \gamma^{T(t)-t-1} \rho_{t+1:T(t)-1} \bar{G}_{t:T(t)}\right)}{\sum_{t \in \mathcal{T}(s,a)} \left((1-\gamma) + (1-\gamma)\sum_{h=t+2}^{T(t)-1} \gamma^{h-t-1} \rho_{t+1:h-1} + \gamma^{T(t)-t-1} \rho_{t+1:T(t)-1}\right)}. $}
    \]

    \textbf{Рекуррентные соотношения.}
    Обозначим числитель и знаменатель добавки от момента $t$ как $N_t$ и $D_t$. Используя $\bar{G}_{t:h} = R_{t+1} + \bar{G}_{t+1:h}$, получаем:
    \[
    \begin{aligned}
    N_t &= R_{t+1} \cdot D_t + \gamma \rho_{t+1} N_{t+1}, \\
    D_t &= (1-\gamma) + \gamma \rho_{t+1} D_{t+1},
    \end{aligned}
    \]
    где
    \[
    \rho_{t+1} = \frac{\pi(A_{t+1}\mid S_{t+1})}{b(A_{t+1}\mid S_{t+1})},
    \qquad N_{T-1} = R_T, \qquad D_{T-1} = 1.
    \]

    \noindent\textbf{Алгоритм.}
    \par\noindent
    \begin{algorithm}[H]
        \DontPrintSemicolon
        \SetAlgoLined
        \KwIn{Поведенческая политика $b$ с покрытием $\pi$}
        \KwOut{Целевая политика $\pi \approx \pi_*$}
        \BlankLine
        \textbf{Инициализация:}\;
        \For{всех $s \in \mathcal{S}$, $a \in \mathcal{A}(s)$}{
            $Q(s,a) \leftarrow$ произвольно\;
            $C_{\text{num}}(s,a) \leftarrow 0$, $\quad C_{\text{den}}(s,a) \leftarrow 0$\;
        }
        $\pi(s) \leftarrow \arg\max_a Q(s,a)$ для всех $s \in \mathcal{S}$\;
        \BlankLine
        \ForEach{эпизод}{
            Сгенерировать эпизод по $b$: $S_0, A_0, R_1, \ldots, S_{T-1}, A_{T-1}, R_T$\;
            $N \leftarrow 0$, $\quad D \leftarrow 0$\;
            \For{$t = T-1, T-2, \ldots, 0$}{
                \eIf{$t = T-1$}{
                    $N \leftarrow R_{t+1}$, $\quad D \leftarrow 1$\;
                }{
                    $\rho \leftarrow \dfrac{\pi(A_{t+1}\mid S_{t+1})}{b(A_{t+1}\mid S_{t+1})}$\;
                    $N \leftarrow R_{t+1} \cdot D + \gamma \cdot \rho \cdot N$\;
                    $D \leftarrow (1-\gamma) + \gamma \cdot \rho \cdot D$\;
                }
                $C_{\text{num}}(S_t, A_t) \leftarrow C_{\text{num}}(S_t, A_t) + N$\;
                $C_{\text{den}}(S_t, A_t) \leftarrow C_{\text{den}}(S_t, A_t) + D$\;
                $Q(S_t, A_t) \leftarrow C_{\text{num}}(S_t, A_t) / C_{\text{den}}(S_t, A_t)$\;
                $\pi(S_t) \leftarrow \arg\max_a Q(S_t, a)$\;
                \If{$A_t \neq \pi(S_t)$}{
                    \textbf{выйти из внутреннего цикла}\;
                }
            }
        }
    \end{algorithm}
\end{solution}
